\section{LastStraw}
\label{sec:benchmark}

% Figure 2: how tasks, human labels and agent runs become frustration verdicts.
\begin{figure}[t]
\centering
\begin{tikzpicture}[
    font=\footnotesize,
    stage/.style={draw=scaleMediumGray, fill=scalePanel, rounded corners=3pt, align=center,
                  text width=2.85cm, minimum height=1.75cm, inner sep=4pt},
    final/.style={draw=scaleCoral, fill=scaleCoral!8, rounded corners=3pt, align=center,
                text width=2.85cm, minimum height=1.75cm, inner sep=4pt},
    arr/.style={-{Latex[length=2mm]}, thick, scaleGray}]
  \node[stage] (author) {\textbf{1. Authoring}\\ \tbd{75} tasks in \tbd{8} categories, a checker per task};
  \node[stage, right=0.3cm of author] (effort) {\textbf{2. Human effort}\\ \tbd{n} annotators do each task: time $t^{H}$, actions $a^{H}$};
  \node[stage, right=0.3cm of effort] (tol) {\textbf{3. Human tolerance}\\ wait budget $\tau$ and action budget $\kappa$ per task};
  \node[stage, right=0.3cm of tol] (runs) {\textbf{4. Agent runs}\\ \tbd{M} models $\times$ \tbd{k} runs, timed screen and action logs};
  \node[final, right=0.3cm of runs] (score) {\textbf{5. Scoring}\\ \tsr, \frate, tolerance gap, pattern labels};
  \draw[arr] (author) -- (effort);
  \draw[arr] (effort) -- (tol);
  \draw[arr] (tol) -- (runs);
  \draw[arr] (runs) -- (score);
\end{tikzpicture}
\caption{Construction and evaluation pipeline of \bench. Human labels (stages 2--3) set a per-task budget, and a run passes only if it reaches the correct end state within that budget (\S\ref{sec:metrics}). \todo{Replace with the final pipeline figure.}}
\label{fig:pipeline}
\end{figure}


\bench pairs each task with two kinds of human labels: how much effort the task takes a person, and how much time and how many actions a person would give an agent before stopping it and doing the task themselves (Figure~\ref{fig:pipeline}). We describe the design principles (\S\ref{sec:principles}), the tasks (\S\ref{sec:tasks}), the human labels (\S\ref{sec:labels}), the metrics (\S\ref{sec:metrics}) and the frustration patterns (\S\ref{sec:patterns}).

\subsection{Design principles}
\label{sec:principles}

\paragraph{Easy for people.} Every task is one that a person familiar with the application completes quickly and reliably from the task text alone. The median human completion time is \tbd{XX} seconds, and \tbd{XX}\% of tasks take under \tbd{XX} seconds (\S\ref{sec:labels}).

\paragraph{Frustration-prone for agents.} Each task centres on an interaction that is routine for people but where agents plausibly stall, loop or take a detour, such as a drag, a hover menu, a custom widget or an interrupting dialog (Table~\ref{tab:categories}).

\paragraph{Correctness is necessary but not sufficient.} Every task has a deterministic checker for the end state, as in existing benchmarks \citep{xie2024osworld,zhou2024webarena}. A correct run that exceeds the task's human budget still fails.

\paragraph{The user's view of time.} Budgets are stated in wall-clock time and in actions, because a user waiting for an agent perceives both how long it takes and how much it flails. \todo{Decide whether model latency counts fully toward wall-clock time or whether we also report a latency-normalized variant; see \S\ref{sec:limitations}.}

\subsection{Tasks}
\label{sec:tasks}

% Table 1: task categories. Category names and examples are a proposed grouping; counts and
% human statistics are pending.
\begin{table}[t]
\centering\small\setlength{\tabcolsep}{5pt}
\resizebox{\textwidth}{!}{%
\begin{tabular}{>{\raggedright\arraybackslash}p{3.6cm}>{\raggedright\arraybackslash}p{6.8cm}rrr}
\toprule
Category & Example task & Tasks & Human time (s) & Tolerance $\tau$ (s) \\
\midrule
Continuous manipulation & Drag the invoice file into the \emph{2026} folder; set the brightness slider to 40\% & \tbd{XX} & \tbd{XX} & \tbd{XX} \\
Scrolling and hidden content & Open the last entry in the sidebar's nested list & \tbd{XX} & \tbd{XX} & \tbd{XX} \\
Transient UI & Choose \emph{Duplicate} from the hover menu of the third row & \tbd{XX} & \tbd{XX} & \tbd{XX} \\
Custom widgets & Set the start date in the date picker to the first Monday of next month & \tbd{XX} & \tbd{XX} & \tbd{XX} \\
Text selection and editing & Make only the second sentence of the paragraph bold & \tbd{XX} & \tbd{XX} & \tbd{XX} \\
Interruptions & Close the newsletter pop-up, then add the item to the cart & \tbd{XX} & \tbd{XX} & \tbd{XX} \\
Keyboard and windows & Copy the order number from the e-mail into the refund form & \tbd{XX} & \tbd{XX} & \tbd{XX} \\
Visual and spatial & Rotate the shape on the canvas so its arrow points down & \tbd{XX} & \tbd{XX} & \tbd{XX} \\
\midrule
\textbf{Total} & & \textbf{75} & \tbd{XX} & \tbd{XX} \\
\bottomrule
\end{tabular}}
\caption{Task categories in \bench with an example task, the number of tasks, and the median human completion time and median tolerance across annotators. \todo{Confirm the category scheme and examples against the 75 tasks.}}
\label{tab:categories}
\end{table}


\paragraph{Authoring.} \tbd{N} annotators wrote the 75 tasks following written guidelines. A task must (i)~have a single unambiguous goal stated in one or two sentences, (ii)~take a person under \tbd{XX} seconds, (iii)~have an end state that a script can verify, and (iv)~avoid logins, payments, CAPTCHAs and actions on live third-party services. \todo{Describe the authoring guidelines and review rounds in Appendix~\ref{app:authoring}.}

\paragraph{Categories.} We group tasks by the interaction that makes them frustration-prone (Table~\ref{tab:categories}): continuous manipulation, scrolling and hidden content, transient UI, custom widgets, text selection and editing, interruptions, keyboard and window management, and visual and spatial tasks. \todo{Check that each task maps to one primary category.}

\paragraph{Environment.} Tasks run in \tbd{Ubuntu 22.04} virtual machines at a screen resolution of \tbd{1920$\times$1080}, with \tbd{Chrome, LibreOffice, a file manager and locally hosted web apps}. Each task ships with a setup script that restores a snapshot and a checker that inspects the final application or file state. \todo{Environment details, application versions, and whether we build on an existing harness such as OSWorld.}

\paragraph{Selection.} \todo{State whether tasks were filtered against any model during authoring (for example, keeping only tasks on which a pilot agent struggled). If so, name the pilot model and report its results separately to avoid inflating its frustration rate.}

\subsection{Human effort and tolerance labels}
\label{sec:labels}

\paragraph{Effort.} Each task was performed by \tbd{n} annotators who did not write it. We record the time from the first view of the task to completion and the number of input actions (clicks, drags, scrolls, key presses and typed strings). For task $i$, the human effort is the median completion time $t^{H}_i$ and the median action count $a^{H}_i$ across annotators. Annotators completed \tbd{XX}\% of attempts on the first try.

\paragraph{Tolerance.} After completing a task, annotators answered two questions: \emph{``If you had asked an assistant to do this task while you watched the screen, how long would you wait before stopping it and doing it yourself?''} and \emph{``How many actions that do not move toward the goal would you accept before stopping it?''} The task's time tolerance $\tau_i$ and action tolerance $\kappa_i$ are the medians of these answers. Across tasks, the median $\tau_i$ is \tbd{XX} seconds, or \tbd{X.X}$\times$ the human completion time. \todo{Decide how wasted actions are counted against $\kappa_i$, or switch $\kappa_i$ to a total-action budget.}

\paragraph{Validating stated tolerance.} Stated tolerance can differ from behaviour. On \tbd{XX} tasks, \tbd{n} annotators watched replays of agent runs and pressed a button at the moment they would take over. \todo{Report how behavioural take-over times compare with the stated $\tau_i$, and inter-annotator agreement on $\tau_i$.}

\subsection{Metrics}
\label{sec:metrics}

For a run $r$ on task $i$, let $c_r \in \{0,1\}$ be the checker's verdict on the final state, $T_r$ the wall-clock time from the instruction to the agent's final action or completion message, and $A_r$ the number of actions it took. Runs are capped at \tbd{XX} actions or \tbd{XX} minutes, well above every task's budget, so that we can also measure correctness without a budget.

\paragraph{Success rate (\sr).} The fraction of runs with $c_r = 1$, as reported by existing benchmarks.

\paragraph{Tolerated success rate (\tsr).} A run succeeds only if it is correct and within both budgets:
\begin{equation}
s_r = c_r \cdot \mathbb{1}[T_r \le \tau_i] \cdot \mathbb{1}[A_r \le \kappa_i], \qquad \tsr = \frac{1}{|R|}\sum_{r \in R} s_r .
\end{equation}
The \emph{frustration rate} is $\frate = 1 - \tsr$. Every run that fails, stops early, or runs past the point at which the user would have taken over counts as a frustration.

\paragraph{Tolerance gap.} $\Delta = \sr - \tsr$ is the share of runs that a correctness-only benchmark counts as successes but a user would not.

\paragraph{Effort ratios.} On correct runs we report the median time ratio $T_r / t^{H}_i$ and action ratio $A_r / a^{H}_i$, which show how far agents are from human effort even when they succeed.

\paragraph{Consistency.} A user who is frustrated once may stop delegating the task. Following pass$^k$ \citep{yao2025taubench}, $\tsr^{k}$ is the fraction of tasks on which all $k$ runs are tolerated successes.

\subsection{Frustration patterns}
\label{sec:patterns}

% Table 2: taxonomy of frustration patterns and how each is detected.
\begin{table}[t]
\centering\small\setlength{\tabcolsep}{5pt}
\resizebox{\textwidth}{!}{%
\begin{tabular}{>{\raggedright\arraybackslash}p{2.9cm}>{\raggedright\arraybackslash}p{8.4cm}>{\raggedright\arraybackslash}p{4.6cm}}
\toprule
Pattern & What the user sees & Detector \\
\midrule
Looping & The agent repeats an action, or a short cycle of actions, without the screen changing in a way that brings it closer to the goal. & Repeated action $n$-grams on near-identical screenshots \\
Stalling & Long pauses before trivial actions, so the screen stays still while the agent deliberates. & Inter-action gap $>\,$\tbd{X}\,s \\
Detour & The agent reaches the goal by a much longer path than needed, e.g.\ menus and search instead of a visible control. & $A > \tbd{X}\times a^{H}$ on correct runs \\
Primitive failure & The agent cannot perform a basic interaction such as a drag, a hover or a scroll inside a panel, and retries it. & LLM judge, human-verified \\
Grounding miss & Clicks or types into the wrong element, sometimes repeatedly. & LLM judge, human-verified \\
False completion & The agent reports success although the end state is wrong. & Checker fails after a ``done'' claim \\
Unneeded hand-back & The agent stops to ask the user for confirmation or input that the task does not require, or gives up. & LLM judge, human-verified \\
Collateral change & The agent completes the task but changes something else that the user must undo. & State diff outside the task scope \\
\bottomrule
\end{tabular}}
\caption{Frustration patterns. A run can show several patterns; patterns explain why a run fails, while the pass criterion itself depends only on correctness and the tolerance budget (\S\ref{sec:metrics}). \todo{Finalize detectors and thresholds; report agreement with human labels in Appendix~\ref{app:detectors}.}}
\label{tab:patterns}
\end{table}


The metrics say whether a run frustrates, and the patterns in Table~\ref{tab:patterns} say how. Looping, stalling and detours are detected from the action log and screenshots with fixed rules; primitive failures, grounding misses and unneeded hand-backs are labelled by an LLM judge whose labels annotators verify on \tbd{XX} runs; false completions and collateral changes come from the checker and a state diff. \todo{Report detector precision and recall against human labels (Appendix~\ref{app:detectors}).}
